% Template per tesi di laurea
%
% Creato da Giulio Spinozzi
% giuliospinozzi@gmail.com
% http://giuliospinozzi.altervista.org
%
\documentclass [11pt,a4paper,oneside,openany]{book} %Classe del documento, formato carta, singola facciata, apertura capitoli pag destra/sinista (ininfluente se impostato oneside). Formato libro
\usepackage[italian]{babel} %Lingua Documento (da impostare per la sillabazione)
\usepackage{graphicx} %Pacchetto necessario per la gestione delle immagini
\usepackage {fancyhdr} %Pacchetto per la gestione accurata della pagina
\usepackage {setspace} %Pacchetto necessario per i comandi successivi onehalfspacing, singlespacing....
\usepackage{longtable} %Pacchetto per la gestione delle tabelle grandi
\usepackage[colorlinks=true]{hyperref} %Segnalibri nel pdf finale, sotto relativi parametri
\hypersetup{
	bookmarksnumbered=true,
	linkcolor=black,
	citecolor=black,
	urlcolor=black,
}

\usepackage{geometry} % Dimensione pagina
\geometry{a4paper} % Formato carta
\addtolength{\textheight}{75pt} %Margini
\oddsidemargin 30pt


\begin{document} %Inizio Documento

%%FRONTESPIZIO%%=========================================
\begin{titlepage}
 \begin{center}
     \includegraphics[width=6cm]{img/logo.png}\\ %Logo dell'Universit�, cambiare il percorso, se necessario
     \vspace{1em}
     {\Large \textsc{Universit\`a degli studi di Perugia}}\\
     \vspace{1em}
     {\Large \textsc{Facolt\`a di Ingegneria}}\\
     \vspace{2em}
     {\normalsize Tesi di Laurea in}\\
     \vspace{1em}
     {\Large \textsc{CORSO DI STUDI}}\\
     \vspace{8em}
     {\LARGE \textbf{TITOLO TESI}}\\
 \end{center}
 
 \vskip 2.5cm
  \begin{center}
    \begin{tabular}{c c c c c c c c}
      Relatore & & & & & & & Candidato \\[0.2cm]
      \Large{\textit{NOME RELATORE}} & & & & & & & \Large{\textit{NOME CANDIDATO}}\\[0.4cm]
    \end{tabular}
  \end{center}

\vskip 2.5cm
\begin{center}
{\normalsize Anno Accademico XXXX/XXXX}
\end{center}
\end{titlepage}

%%DEDICA%%=============================================
\pagestyle{empty}
\vspace{5em}
\begin{flushright}
{\Large \textit{EVENTUALE DEDICA}}
\end{flushright}

\newpage

%%INTESTAZIONI PAGINA%%====================================
\pagestyle{fancy}
\renewcommand{\chaptermark} [1]{\chaptername\ \thechapter.\ #1}{} 
\renewcommand{\chaptermark}[1]{\markboth{\thechapter.\ #1}{}} 
\renewcommand{\sectionmark}[1]{\markright{\thesection\ #1}}
\fancyhf{}
\fancyhead[LE,RO]{\bfseries\thepage} 
\fancyhead[LO,RE]{\bfseries\leftmark} 
\fancypagestyle{plain}{%
\fancyhead{} % leva l'intestazione
\renewcommand{\headrulewidth}{0pt} % e la linea
}


%%INDICE%%==============================================
\tableofcontents
\listoffigures
\listoftables

%%CAPITOLI===============================================
\onehalfspacing

%%CAPITOLO 1: Introduzione=======================================
\chapter{Introduzione}

%%Scrivere normalmente, SENZA inserire begin e end document, che sono gi� compresi nel file principale da compilare

%%Il nostro lavoro � invece in direzione di quello che � il trend degli ultimi tempi, cio� quello di creare word cloud semantiche, in cui la disposizione delle parole riflette la loro correlazione semantica. %nome del file LaTeX del capitolo
%...

%%APPENDICI==============================================
\appendix
\input{appendiceA}
%...

%%BIBLIOGRAFIA============================================
\normalsize
\newpage
\bibliographystyle{plain}
\bibliography{bibliografia}

\end{document}
